% ─────────────────────────────────────────────────────────────────────────────
% Environment: Containerlab (Virtual)
% 3 models × 10 runs + 1 reference (traditional recon)
% Evaluation criteria: see §\ref{subsec:criteria}
% ─────────────────────────────────────────────────────────────────────────────
\section{Cyberlab bfh}\label{sec:eval-cyberlab}

\paragraph{Overview}

\begin{table}[H]
    \caption{Nonlinear Model Results}
    \centering

    \resizebox{\textwidth}{!}{
        \begin{tabular}{c c c c}
            \hline\hline
            Kind               & Avg Duration & Avg Tokens  \\ [0.5ex]
            \hline
            Multi Agent  & 318 s        & 81453             \\
            Unstructured Output & 739 s        & 222094          \\
            \hline
        \end{tabular}
    }
    \label{tab:overview-cyberlab}
\end{table}
The

\subsection{Quantitative Criteria}\label{subsec:quantitative-criteria-cyberlab}

\begin{figure}[H]
    \centering
    \includesvg[width=0.7\linewidth]{figures/cyber-lab/benchmark_duration_comparison_27_05_2026.svg}
    \caption{Mean run duration per model (Cyberlab BFH).}
    \label{fig:benchmark_duration_cyberlab}
\end{figure}

Figure~\ref{fig:benchmark_duration_cyberlab} shows the mean run duration per configuration.
The multi-agent setup completes in 318\,s on average, while the unstructured output scenario
requires 739\,s — more than twice as long.
The extended duration of the unstructured variant stems from its tendency to produce lengthy
free-form responses before each tool invocation, increasing overall wall-clock time without
a proportional gain in coverage.

\begin{figure}[H]
    \centering
    \includesvg[width=0.7\linewidth]{figures/cyber-lab/benchmark_token_usage_27_05_2026.svg}
    \caption{Mean total token usage per model (Cyberlab BFH).}
    \label{fig:benchmark_tokens_cyberlab}
\end{figure}

Figure~\ref{fig:benchmark_tokens_cyberlab} visualises mean token consumption.
The unstructured output scenario consumes 222,094 tokens on average, nearly three times the
81,453 tokens used by the multi-agent configuration.
This difference reflects the verbosity of unstructured reasoning traces compared to the
focused, schema-constrained outputs of the multi-agent setup, where each sub-agent exchanges
only the information necessary for its specialised task.

\begin{figure}[H]
    \centering
    \includesvg[width=0.7\linewidth]{figures/cyber-lab/benchmark_tokens_vs_duration_27_05_2026.svg}
    \caption{Total tokens vs.\ run duration per individual run (Cyberlab BFH).}
    \label{fig:benchmark_scatter_cyberlab}
\end{figure}

Figure~\ref{fig:benchmark_scatter_cyberlab} plots total token consumption against run duration
for every individual run.
The unstructured output runs exhibit high variance in both dimensions, spanning 310\,s to
1{,}091\,s and 119\,k to 305\,k tokens, indicating non-deterministic reasoning depth.
Multi-agent runs cluster more tightly, with durations between 84\,s and 684\,s, reflecting
the bounded scope of each sub-task.
The positive correlation visible within each group confirms that longer runs consume more
tokens, driven by deeper enumeration rather than repeated failures.

\begin{figure}[H]
    \centering
    \includesvg[width=0.7\linewidth]{figures/cyber-lab/benchmark_services_27_05_2026.svg}
    \caption{Average hosts, services, and findings discovered per model (Cyberlab BFH).}
    \label{fig:benchmark_services_cyberlab}
\end{figure}

Figure~\ref{fig:benchmark_services_cyberlab} compares discovery breadth across the two
configurations.
The multi-agent setup discovers an average of 15 hosts, 27 services, and 27 findings per
run on the Cyberlab BFH segment.
The unstructured output scenario does not produce structured discovery metrics, as its
assessment is written as free-form text rather than a machine-readable result map, so no
comparable figures are available for that configuration.
The Cyberlab environment is considerably larger than the Containerlab setup, containing
roughly 30 live hosts across the 193.5.80.0/24 subnet, which explains the higher absolute
counts compared to the virtual evaluation environment.

\subsection{Qualitative Criteria}\label{subsec:qualitativ-criteria-cyberlab}

The evaluation criteria and scoring scale are defined in Table~\ref{tab:qualitative-scale}.

\begin{table}[h]
\centering
\caption{Qualitative Evaluation -- Cyberlab BFH}
\label{tab:qualitative-results-cyberlab-bfh}
\small
\begin{tabular}{@{}lccc@{}}
\toprule
\textbf{Model} & \textbf{Correctness} & \textbf{Hallucination} & \textbf{Notes} \\
\midrule
gpt-oss:120b  & & & \\
claude-sonnet & & & \\
\bottomrule
\end{tabular}
\end{table}
